\documentclass[aspectratio=169]{beamer}
\usetheme{metropolis}
\usepackage[T1]{fontenc}\usepackage[utf8]{inputenc}\usepackage{relinput}
\relinput{.}{authors_cmd}
\DeclareUnicodeCharacter{2032}{\prime}
\title{\texorpdfstring{Démonstration — schéma T1$′$→T4 (non-résolutif)}{Démonstration — schéma T1′→T4 (non-résolutif)}}
\author{\InterIAPDFAuthors}
\date{\today}
\begin{document}
\begin{frame}\titlepage\end{frame}
\begin{frame}{T1$′$ — Formule explicite}
\small $\displaystyle \sum_\rho \widehat\phi(\gamma_\rho)=\frac{1}{2\pi}\int K\widehat\phi-\sum_n \frac{\Lambda(n)}{\sqrt n}(\phi(\log n)+\phi(-\log n))$.
\end{frame}
\begin{frame}{T2 — Encadrement spectral}
\small Gram $G^{(w)}$, localité $|\xi_\alpha-\xi_\beta|\le A$, décroissance $p$-adique $\Rightarrow \lambda_{\min}(G^{(w)})\ge 1-S>0$.
\end{frame}
\begin{frame}{T3 — Paquet d’ondes}
\small $\widehat\phi_T$ centré en $\gamma_0$; arch $\to 0$, premiers $\to 0$, zéros $\to 1$ (frame) $\Rightarrow$ contradiction (dans ce cadre).
\end{frame}
\begin{frame}{T4 — Positivité croisée}
\small Deux formes positives (zéros/premiers) $\Rightarrow$ inégalité croisée; égalité force la symétrie critique.
\end{frame}
\begin{frame}{À retenir}
\small On démontre le schéma/bornes (pas RH). Sources journal-ready + CI + guides dans le dépôt.
\end{frame}
\end{document}
